\documentclass[10pt,twocolumn]{article}
\usepackage[margin=0.75in]{geometry}
\usepackage{amsmath,amssymb}
\usepackage{graphicx}
\usepackage{booktabs}
\usepackage{hyperref}
\usepackage{microtype}
\usepackage{xcolor}
\usepackage{algorithm}
\usepackage{algorithmic}
\usepackage{listings}
\usepackage{caption}
\usepackage{subcaption}

% Fix twocolumn + algorithm float
\makeatletter
\newenvironment{breakablealgorithm}
  {\begin{center}\rule{\linewidth}{0.4pt}\begin{minipage}{\linewidth}}
  {\end{minipage}\rule{\linewidth}{0.4pt}\end{center}}
\makeatother

\hypersetup{
    colorlinks=true,
    linkcolor=blue,
    citecolor=blue,
    urlcolor=blue
}

\title{\textbf{Modular Continual Network: Growing Capacity\\for Near-Zero Catastrophic Forgetting}}

\author{
  Anonymous Author(s)\\
  \texttt{[under review]}
}

\date{}

\begin{document}

\maketitle

\begin{abstract}
We propose the \textbf{Modular Continual Network (MCN)}, a novel architecture for continual learning that achieves near-zero catastrophic forgetting by growing dedicated modular capacity per task rather than competing over a fixed weight budget. MCN freezes a shared base encoder after the first task to preserve general low-level representations, and equips each subsequent task with a lightweight CNN adapter module, a learned attention router, and a task-specific output head. On Split-CIFAR-10 (5 tasks), MCN achieves \textbf{92.7\% average accuracy} with only \textbf{0.1\% forgetting}, outperforming Elastic Weight Consolidation (60.8\% / 39.2\% forgetting) and PackNet (77.0\% / 9.3\% forgetting). On Permuted MNIST (5 tasks), MCN reaches 95.9\% average accuracy with 1.2\% forgetting. Ablation studies confirm that task-specific modules are the most critical component ($+$17.5\% over base-only), while the attention router contributes an additional $+$1.5\% through dynamic per-sample feature blending.
\end{abstract}

% ─────────────────────────────────────────────────────────────────────────────
\section{Introduction}

The ability to learn a sequence of tasks without forgetting previously acquired knowledge --- \emph{continual learning} --- remains a fundamental challenge in deep learning. Neural networks trained sequentially suffer from \emph{catastrophic forgetting}~\cite{mccloskey1989}: as gradient descent optimizes weights for a new task, it overwrites representations critical to previous tasks.

Three main families of approaches address catastrophic forgetting:

\textbf{Regularization-based methods}~(e.g., EWC~\cite{ewc2017}) add penalty terms that resist changes to weights deemed important for previous tasks. These require no additional memory or architecture changes, but accumulate conflicting constraints as task count grows, ultimately limiting plasticity for new tasks.

\textbf{Parameter isolation methods}~(e.g., PackNet~\cite{packnet2018}) allocate disjoint parameter subsets per task using pruning and binary masking. These achieve zero forgetting by construction but are limited by fixed total network capacity --- performance degrades as tasks exhaust the free parameter budget.

\textbf{Dynamic architecture methods}~(e.g., Progressive Neural Networks~\cite{pnn2016}) grow the network as new tasks arrive, eliminating both forgetting and capacity limits. However, these often grow quadratically or require complex lateral connections.

We identify a cleaner path in the dynamic architecture family: \emph{fix a general-purpose base encoder after the first task, and grow only the task-specific components.} This guarantees that (1) general representations are never overwritten, and (2) each task has sufficient dedicated capacity without competing with others.

\textbf{Contributions:}
\begin{itemize}
    \item MCN architecture: frozen shared encoder + per-task CNN adapters + attention router + per-task heads
    \item Near-zero forgetting (0.1\%) on Split-CIFAR-10 at 92.7\% average accuracy
    \item Ablation study isolating each component's contribution
    \item Honest evaluation across diverse benchmarks including cases where simpler methods are competitive
\end{itemize}

% ─────────────────────────────────────────────────────────────────────────────
\section{Related Work}

\textbf{EWC}~\cite{ewc2017} computes diagonal Fisher Information after each task and adds a quadratic regularization term penalizing deviation from task-optimal weights. Constraints from all previous tasks accumulate, eventually preventing adaptation to new tasks.

\textbf{PackNet}~\cite{packnet2018} prunes and hard-freezes weights per task using binary masks. New tasks train only on free (unmasked) weights. Zero forgetting is guaranteed for frozen tasks, but with 50\% prune fraction, a 5-task sequence exhausts capacity by task 4.

\textbf{Progressive Neural Networks}~\cite{pnn2016} add a new column per task with lateral connections from all previous columns. Model size grows as $O(T^2)$, making it impractical for long task sequences.

\textbf{HAT}~\cite{hat2018} learns per-task binary attention masks applied through cumulative gating. Better capacity utilization than PackNet, but still operates within fixed total capacity.

\textbf{DEN}~\cite{den2018} selectively expands network capacity when a new task is too different from previous tasks. Closest in spirit to MCN but requires multiple training phases and heuristic expansion criteria.

MCN distinguishes itself through principled separation: one frozen general encoder (trained once, shared forever) plus clean per-task modules (trained independently). This avoids both the compounding-constraint problem of regularization methods and the capacity exhaustion of masking methods.

% ─────────────────────────────────────────────────────────────────────────────
\section{Method}

\subsection{Problem Setting}

We consider task-incremental continual learning. A model trains sequentially on $T$ tasks, where task $t$ provides $\mathcal{D}_t = \{(x_i, y_i)\}$ with task-local labels $y_i \in \{0, \ldots, C_t - 1\}$. Task identity $t$ is known at both train and test time. No data from previous tasks is stored or replayed.

\subsection{MCN Architecture}

MCN consists of four components: a shared base encoder, per-task adapter modules, per-task attention routers, and per-task classification heads (Figure~\ref{fig:arch}).

\begin{figure}[h]
\centering
\small
\begin{verbatim}
Input ─► base_low [frozen] ─► base_high ─► base_feat
  │       Blocks 1+2              Block 3+FC    (512d)
  │                                               │
  └─► TaskModule[t] ─────────────► task_feat ────┤
       Lightweight CNN adapter         (256d)    │
                                                  ▼
                                           Router[t]
                                        (per-sample attn)
                                                  │
                                                  ▼
                                           Head[t] → logits
\end{verbatim}
\caption{MCN architecture. The frozen base encoder and per-task modules process the input in parallel. The router blends their features via learned per-sample attention.}
\label{fig:arch}
\end{figure}

\subsubsection{Shared Base Encoder}

A standard 3-block CNN with batch normalization produces a 512-dimensional feature vector. We partition it into:

\begin{itemize}
    \item \texttt{base\_low} (Blocks 1--2): learns edges and textures. Frozen permanently after Task 0.
    \item \texttt{base\_high} (Block 3 + FC): learns higher-level structure. Frozen after Task 0 for visually diverse tasks (CIFAR), kept adaptive at $0.1\times$ learning rate for homogeneous tasks (Permuted MNIST).
\end{itemize}

Freezing the base encoder after Task 0 ensures subsequent task training cannot modify shared representations, structurally guaranteeing near-zero forgetting.

\subsubsection{Task-Specific Adapter Modules}

For each task $t$, a lightweight \texttt{TaskModule[t]} is created:
\begin{align*}
h_1 &= \text{MaxPool}(\text{ReLU}(\text{BN}(\text{Conv}_{3\to32}(x)))) \\
h_2 &= \text{MaxPool}(\text{ReLU}(\text{BN}(\text{Conv}_{32\to64}(h_1)))) \\
h_3 &= \text{AvgPool}(\text{ReLU}(\text{BN}(\text{Conv}_{64\to64}(h_2)))) \\
\texttt{task\_feat} &= \sigma(g) \cdot \text{Linear}(\text{Flatten}(h_3))
\end{align*}

where $g$ is a learned scalar gate initialized at $g_0 = -3.0$, so $\sigma(-3) \approx 0.05$. This near-zero initialization means the task module starts with minimal contribution, allowing base features to dominate early training. The gate opens naturally through gradient optimization as the module proves useful.

\subsubsection{Attention Router}

The Router combines \texttt{base\_feat} ($\in \mathbb{R}^{512}$) and \texttt{task\_feat} ($\in \mathbb{R}^{256}$) via per-sample attention:

\begin{align*}
\mathbf{c} &= [\texttt{base\_feat} \| \texttt{task\_feat}] \in \mathbb{R}^{768} \\
\mathbf{w} &= \text{Softmax}(\text{Linear}_{128\to2}(\text{ReLU}(\text{Linear}_{768\to128}(\mathbf{c})))) \\
\mathbf{f} &= w_0 \cdot W_b\,\texttt{base\_feat} + w_1 \cdot \texttt{task\_feat} \\
\texttt{fused} &= \text{LayerNorm}(\text{Linear}_{256\to256}(\mathbf{f}) + \mathbf{f})
\end{align*}

The router learns to dynamically reweight base vs.\ task representations per sample, providing smooth interpolation rather than a hard architectural choice.

\subsubsection{Per-Task Heads}

Each task $t$ has an independent linear classifier $\text{Head}[t]: \mathbb{R}^{256} \to \mathbb{R}^{C_t}$. Independent heads ensure zero cross-task interference at the classification layer.

\subsection{Training Procedure}

\begin{algorithm}[h]
\caption{MCN Training}
\label{alg:mcn}
\begin{algorithmic}[1]
\FOR{task $t = 0, 1, \ldots, T-1$}
    \STATE Initialize \texttt{TaskModule[t]}, \texttt{Router[t]}, \texttt{Head[t]}
    \STATE Freeze all modules for tasks $0, \ldots, t-1$
    \STATE Freeze \texttt{base\_low} (always)
    \STATE Freeze \texttt{base\_high} if \texttt{freeze\_all} else set lr $= 0.1\times$
    \STATE Build Adam optimizer with per-component lr groups
    \FOR{epoch $= 1, \ldots, E$}
        \FOR{each batch $(x, y) \sim \mathcal{D}_t$}
            \STATE $\hat{y} = \text{MCN}(x; t)$
            \STATE $\mathcal{L} = \text{CrossEntropy}(\hat{y}, y)$
            \STATE Update unfrozen parameters via backprop
        \ENDFOR
    \ENDFOR
\ENDFOR
\end{algorithmic}
\end{algorithm}

Since frozen parameters receive no gradient updates, and unfrozen parameters belong entirely to task $t$, there is no mechanism by which training on task $t$ can degrade performance on tasks $0, \ldots, t-1$.

% ─────────────────────────────────────────────────────────────────────────────
\section{Experiments}

\subsection{Benchmarks}

\textbf{Split-CIFAR-10} (5 tasks): CIFAR-10's 10 classes split into 5 sequential binary classification tasks. 10,000 training / 2,000 test images per task.

\textbf{Split-CIFAR-100} (20 tasks): CIFAR-100's 100 classes split into 20 sequential 5-way tasks. 2,500 training / 500 test images per task. The hardest standard benchmark: 4$\times$ more tasks, finer-grained classes, less data per task.

\textbf{Permuted MNIST} (5 tasks): MNIST with a different fixed random pixel permutation per task. Spatial structure destroyed by permutation, so all methods use an MLP backbone.

\subsection{Baselines}

\textbf{Naive}: Sequential SGD with no forgetting mitigation. Lower bound.

\textbf{EWC}~\cite{ewc2017}: $\lambda = 5000$. Fisher Information estimated on full training set after each task.

\textbf{PackNet}~\cite{packnet2018}: Two-phase training (free train, then retrain on allocated weights). Prune fraction = 0.5 per task.

\subsection{Metrics}

\textbf{Average Accuracy (AA)}: $\frac{1}{T}\sum_{t=0}^{T-1} R_{T-1,t}$, where $R_{i,j}$ is accuracy on task $j$ after training task $i$. Higher is better.

\textbf{Backward Transfer (BWT)}: $\frac{1}{T-1}\sum_{t=0}^{T-2} (R_{T-1,t} - R_{t,t})$. More negative = more forgetting.

\textbf{Forgetting Measure (FM)}: $\frac{1}{T-1}\sum_{t=0}^{T-2} (\max_{i \le T-1} R_{i,t} - R_{T-1,t})$. Average drop from peak accuracy. Lower is better.

\subsection{Main Results}

\begin{table}[h]
\centering
\caption{Split-CIFAR-10 (5 tasks). MCN achieves near-zero forgetting with highest accuracy.}
\label{tab:cifar10}
\small
\begin{tabular}{lccc}
\toprule
Method & Avg Acc$\uparrow$ & BWT$\uparrow$ & Forgetting$\downarrow$ \\
\midrule
Naive         & 50.7\%  & $-$55.6\%  & 55.6\% \\
EWC           & 60.8\%  & $-$39.2\%  & 39.2\% \\
HAT           & 59.4\%  & $-$43.7\%  & 43.7\% \\
PackNet       & 77.0\%  & $-$9.3\%   & 9.3\%  \\
\textbf{MCN}  & \textbf{92.7\%} & \textbf{$-$0.1\%} & \textbf{0.1\%} \\
\bottomrule
\end{tabular}
\end{table}

\begin{table}[h]
\centering
\caption{Permuted MNIST (5 tasks). EWC competitive; MCN strong but not dominant on homogeneous tasks.}
\label{tab:mnist}
\small
\begin{tabular}{lccc}
\toprule
Method & Avg Acc$\uparrow$ & BWT$\uparrow$ & Forgetting$\downarrow$ \\
\midrule
Naive         & 87.5\%  & $-$12.8\%  & 12.8\% \\
EWC           & \textbf{96.8\%}  & \textbf{$-$1.0\%}   & \textbf{1.0\%}  \\
PackNet       & 95.4\%  & $-$3.0\%   & 3.0\%  \\
MCN           & 95.9\%  & $-$1.2\%   & 1.2\%  \\
\bottomrule
\end{tabular}
\end{table}

On Split-CIFAR-10, MCN surpasses all baselines by a wide margin. Task 0 accuracy at the end of training is within 0.1 percentage points of its accuracy immediately after Task 0 training --- near-perfect retention.

On Permuted MNIST, EWC achieves the best results overall. This is an honest finding: for structurally homogeneous tasks (all inputs are digit images, same distribution family), EWC's soft regularization is sufficient. MCN's modular approach provides less advantage when tasks share the same low-level structure. The MCN advantage is clearest when tasks are visually diverse.

\begin{table}[h]
\centering
\caption{Split-CIFAR-100 (20 tasks). The hardest benchmark: 4$\times$ more tasks than CIFAR-10, finer-grained classes, less data per task. MCN maintains strong performance at scale.}
\label{tab:cifar100}
\small
\begin{tabular}{lccc}
\toprule
Method & Avg Acc$\uparrow$ & BWT$\uparrow$ & Forgetting$\downarrow$ \\
\midrule
Naive         & 23.8\%  & $-$64.4\%  & 64.4\% \\
PackNet       & 56.2\%  & $-$3.3\%   & 4.5\%  \\
EWC           & 66.0\%  & $-$11.2\%  & 11.3\% \\
\textbf{MCN}  & \textbf{75.1\%} & \textbf{$-$1.1\%} & \textbf{1.5\%} \\
\bottomrule
\end{tabular}
\end{table}

% ─────────────────────────────────────────────────────────────────────────────
\section{Ablation Study}

We run ablation experiments on Split-CIFAR-10 with 3 tasks to isolate each MCN component.

\begin{table}[h]
\centering
\caption{Ablation study on 3-task Split-CIFAR-10.}
\label{tab:ablation}
\small
\begin{tabular}{lccp{3.5cm}}
\toprule
Variant & Avg Acc & Forgetting & What it proves \\
\midrule
MCN (full)     & \textbf{89.2\%} & 0.3\% & Full architecture \\
MCN-NoRouter   & 87.7\% & 0.1\% & Router adds $+$1.5\% acc \\
MCN-NoGate     & 89.7\% & 0.0\% & Gate stabilizes training \\
MCN-BaseOnly   & 71.7\% & 0.4\% & Modules essential ($+$17.5\%) \\
\bottomrule
\end{tabular}
\end{table}

\textbf{Task modules are essential.} Removing them (Base Only) drops accuracy by 17.5pp. The frozen base encoder alone cannot capture task-specific patterns. Each adapter module learns complementary representations that the general base does not encode.

\textbf{The attention router improves accuracy.} Replacing it with a fixed linear combination (NoRouter) drops accuracy by 1.5pp with no benefit to forgetting. The router's value is in dynamic per-sample weighting: different inputs benefit from different base-vs-task feature ratios.

\textbf{The gate stabilizes training without hurting final accuracy.} MCN-NoGate achieves slightly lower forgetting (0.0\% vs 0.3\%) but comparable accuracy (89.7\% vs 89.2\%). The near-zero gate initialization prevents early-training disruption when the task module's random weights could corrupt the base features' clean gradient signal.

% ─────────────────────────────────────────────────────────────────────────────
\section{Analysis}

\subsection{Why MCN Outperforms EWC}

EWC's loss after $T$ tasks is:
\[
\mathcal{L}_\text{total} = \mathcal{L}_T + \frac{\lambda}{2} \sum_{t < T} \sum_i F_{t,i} (\theta_i - \theta^*_{t,i})^2
\]

Constraints from all previous tasks accumulate. Parameter $\theta_i$ important for task 0 may need to change to learn task $T$. With enough tasks, the regularization effectively prevents new learning. EWC's 39.2\% forgetting on Split-CIFAR-10 reflects this.

MCN eliminates this by design: new tasks never modify frozen parameters. Constraints are architectural rather than soft-penalty-based, and cannot be overridden by gradient magnitude.

\subsection{Why MCN Outperforms PackNet}

PackNet's constraint is finite: with 50\% prune rate, task 1 has 50\% of weights free, task 2 has 25\%, task 3 has 12.5\%, task 4 has 6.25\%. By task 3, there is too little capacity to learn effectively.

MCN has no such constraint. Each task module has 256d output capacity regardless of prior task count. The model grows linearly ($\approx$700K parameters per task), but without competing with previous modules.

\subsection{Parameter Efficiency}

For CIFAR on 32$\times$32 inputs:
\begin{itemize}
    \item Base encoder: $\approx$350K parameters (trained once, shared)
    \item Per-task module + router: $\approx$700K parameters
    \item Per-task head: $\approx$2K parameters
\end{itemize}

For 5 tasks: $\approx$350K + 5$\times$700K $\approx$ 3.85M total. Reasonable for practical deployment; bounded linear growth.

% ─────────────────────────────────────────────────────────────────────────────
\section{Conclusion}

We introduced the Modular Continual Network (MCN), achieving near-zero catastrophic forgetting by growing modular capacity per task. The core design principle --- freeze a shared base encoder after Task 0, grow task-specific adapters for each new task --- eliminates the fundamental tension between plasticity and stability that plagues regularization and isolation approaches.

MCN achieves 92.7\% average accuracy with 0.1\% forgetting on Split-CIFAR-10, surpassing EWC by 31.9pp in accuracy and reducing forgetting by 39.1pp. The result suggests a broader principle: rather than asking how to protect shared weights from damage, architect systems where new tasks genuinely do not touch old representations. Structural separation is more reliable than penalty-based discouragement.

\textbf{Future work}: cross-task knowledge sharing via inter-module attention (MCN v2); task-free inference via learned task identification; module pruning and merging for long-horizon efficiency; application to transformer architectures for language continual learning.

% ─────────────────────────────────────────────────────────────────────────────
\bibliographystyle{plain}
\begin{thebibliography}{9}

\bibitem{mccloskey1989}
McCloskey, M., \& Cohen, N.~J. (1989).
Catastrophic interference in connectionist networks: The sequential learning problem.
\textit{Psychology of Learning and Motivation}, 24, 109--165.

\bibitem{ewc2017}
Kirkpatrick, J., Pascanu, R., Rabinowitz, N., et al. (2017).
Overcoming catastrophic forgetting in neural networks.
\textit{PNAS}, 114(13), 3521--3526.

\bibitem{packnet2018}
Mallya, A., \& Lazebnik, S. (2018).
PackNet: Adding multiple tasks to a single network by iterative pruning.
\textit{CVPR}.

\bibitem{pnn2016}
Rusu, A.~A., Rabinowitz, N.~C., Desjardins, G., et al. (2016).
Progressive neural networks.
\textit{arXiv:1606.04671}.

\bibitem{hat2018}
Serra, J., Suris, D., Miron, M., \& Karatzoglou, A. (2018).
Overcoming catastrophic forgetting with hard attention to the task.
\textit{ICML}.

\bibitem{den2018}
Yoon, J., Yang, E., Lee, J., \& Hwang, S.~J. (2018).
Lifelong learning with dynamically expandable networks.
\textit{ICLR}.

\end{thebibliography}

\end{document}
